\section{Reasoning Categorizations}
To better understand the reasoning capabilities required in \ours, we summarize 4 reasoning types with representative examples from each of the 12 \ours datasets.
\begin{itemize}
    \item Deductive reasoning: The document usually describes a general principle or theorem that could be applied to explain a specific scenario or solve a specific problem present in the query. For example, in Table 20, the general mechanism of meristem is applied to explain the scenario where a tree trunk could sprout and grow after being cut. (See more examples in Table 21, 25, 22, 31).
    \item Analogical reasoning: The document draws parallel with the query in underlying logic, which indicates that the method used in the document could be also used to solve the query problem. For example, in Table 29, although the query problem appears to be different from the one in the documents, they share similar underlying logic and can be both solved by the Chicken McNugget Theorem (See more examples in Table 27, 30).
    \item Causal reasoning: The document and the query present a cause-and-effect relationship, i.e., To fix the problem in the query, we need to find the cause in the document. For example, in Table 24, the query presents a problem where the debug message is missing, which is caused by the parameter settings in the source code (positive documents).
    \item Analytical reasoning: The document provides critical concepts or knowledge that support reasoning chains to solve problems in queries. For example, in Table 26, one will need to first analyze the problem in the query: To determine whether reusing water for plants is beneficial, we need to first understand where the water comes from and what could be contained in the water; As described in the query, the water comes from watering, which goes through plants and soil and finally arrives the plates below the pots; This indicates that the water could carry minerals, soluble salts and other materials in soil and plants. The document provides critical knowledge that soluble salts could be dissolved and accumulate in the water, which will be harmful to plants. This piece of information helps to complete the reasoning chain for deriving the final answer to the query. (See more examples in Table 23, 28).
\end{itemize}

With categorized reasoning capabilities, people who use BRIGHT can now understand better about their achievements if they observe improvements on specific datasets.
We hope that this will facilitate future research on more systematically developing retrieval models with strong reasoning capabilities.